A redefinition of the SI second requires optical clocks in different
laboratories to measure the same frequency ratio independently and to agree
within \refThreshold. Several clocks now report self-evaluated uncertainties
at this level, but no comparison between clocks in different cities has yet
verified the requirement. Here we compare a ytterbium lattice clock in Wuhan
with a strontium lattice clock in Shanghai, about $\siteKm$\,km apart,
through a $\fibreKm$\,km phase-stabilised fibre link, and determine the
Yb/Sr frequency ratio from \nDays~days of coherent operation
(\nSamples~retained seconds across \nSeg~segments). Carrying the
time-varying gravitational-redshift term as an uncertainty item, we obtain
Yb/Sr~=~\Rheadline(\RheadlineUnc), with a statistical uncertainty of
\uStatShort$\times10^{-18}$ and a total standard uncertainty of
\uTotalShort$\times10^{-18}$. The central value differs from the NIST--JILA
result by \dVsNistE and from the European-network result by \dVsEuE, both
within the redefinition threshold. To our knowledge, this is the first
reported inter-city, cross-species frequency-ratio measurement meeting the
uncertainty requirement of the roadmap towards the SI-second redefinition.
